% =====================================================================
%  CHƯƠNG 1 — MỞ ĐẦU VỀ KNOWLEDGE GRAPH VÀ RAG
% =====================================================================
\chapter{Mở đầu về Knowledge Graph và RAG}\label{chap:mo-dau}

Chương này đặt nền cho toàn bài: đồ thị tri thức là gì, RAG là gì, vì sao hai kỹ thuật này bổ sung đúng chỗ yếu
của LLM, và chúng được ghép với nhau như thế nào trong hệ thống của tiểu luận. Hình~\ref{fig:hesinhthai} là bản đồ
của cả báo cáo; mỗi khối trên hình tương ứng một chương phía sau.

\hinhtikz{t1_ecosystem}{Tổng quan hệ sinh thái tri thức đa miền của tiểu luận: nguồn tri thức thô được tổ chức
thành hai kho bổ sung cho nhau là đồ thị tri thức và chỉ mục vector; tác tử ReAct phối hợp hai kho với các mô hình
học sâu để trả lời người dùng}{hesinhthai}

\section{Khái niệm đồ thị tri thức và RAG}

\subsection{Khái niệm đồ thị tri thức (Knowledge Graph, KG)}
Đồ thị tri thức là một đồ thị dữ liệu dùng để tích lũy và truyền đạt tri thức về thế giới thực, trong đó
\textbf{nút} biểu diễn thực thể và \textbf{cạnh} biểu diễn quan hệ giữa các thực thể~\cite{hogan2021kg}. Khác một
cơ sở dữ liệu quan hệ, KG không cố định lược đồ thành các bảng: muốn thêm một loại quan hệ mới chỉ cần thêm một
nhãn cạnh mới. Khác một văn bản, mỗi mẩu tri thức trong KG là một đơn vị nguyên tử, máy đọc được và có thể truy
vấn chính xác.

Về hình thức, gọi $\mathcal{E}$ là tập thực thể, $\mathcal{R}$ là tập quan hệ (vị từ). Một KG là tập các
\textbf{bộ ba} (chủ thể, vị từ, đối tượng):
\begin{equation}
  \mathcal{G} = \{ (s, p, o) \mid s, o \in \mathcal{E},\ p \in \mathcal{R} \} \subseteq \mathcal{E} \times \mathcal{R} \times \mathcal{E}.
  \label{eq:kg}
\end{equation}
Ví dụ, sự kiện ``metformin tương tác nguy hiểm với rượu bia'' trong \texttt{VietnamDiabetesKG} được lưu thành bộ
ba
\[
  (\textsf{Metformin},\ \textsf{TUONG\_TAC\_NGUY\_HIEM},\ \textsf{Ruou\_Bia}).
\]
Đi kèm tập bộ ba là một
\textbf{bản thể học} (ontology) quy định thực thể thuộc lớp nào, vị từ nào được nối những lớp nào; bản thể học
biến một đồ thị nhãn tự do thành tri thức có ngữ nghĩa, để máy kiểm tra được và suy ra được tri thức mới
(Chương~\ref{chap:kg}).

Ba tính chất làm KG có giá trị với hệ thống thông minh:
\begin{itemize}
  \item \textbf{Tường minh và truy vết được.} Mỗi cạnh trong hai KG của bài mang thuộc tính \texttt{source} trỏ về
        tài liệu gốc; câu trả lời dựa trên KG luôn nêu được căn cứ.
  \item \textbf{Đa bước.} Câu hỏi ``Fed tăng lãi suất thì vàng SJC ra sao'' không có câu trả lời trong bất kỳ
        đoạn văn đơn lẻ nào, nhưng là một đường đi qua 3--4 cạnh trên đồ thị nhân quả tài chính.
  \item \textbf{Có ràng buộc.} Bản thể học chặn được bộ ba vô nghĩa như ``rau xanh là thuốc đầu tay của metformin''
        ngay khi nạp dữ liệu, trước khi nó lan vào câu trả lời.
\end{itemize}

\subsection{Khái niệm sinh tăng cường truy xuất (RAG)}
RAG được Lewis và cộng sự đề xuất năm 2020~\cite{lewis2020rag} để giải bài toán tri thức chuyên sâu: thay vì bắt
mô hình nhớ mọi thứ trong tham số, ta cho nó \emph{tra cứu} một kho văn bản ngoài trước khi sinh. Một hệ RAG có
hai thành phần: bộ truy xuất $p_\eta(z \mid x)$ chọn các đoạn văn $z$ liên quan tới câu hỏi $x$, và bộ sinh
$p_\theta(y \mid x, z)$ viết câu trả lời $y$ dựa trên cả câu hỏi lẫn đoạn văn. Xác suất câu trả lời là tổng theo
các đoạn văn được truy xuất:
\begin{equation}
  p(y \mid x) \approx \sum_{z \in \operatorname{top-}k\, p_\eta(\cdot \mid x)} p_\eta(z \mid x)\, p_\theta(y \mid x, z).
  \label{eq:rag}
\end{equation}
Trong thực hành hiện nay, các đoạn top-$k$ được nối vào prompt và LLM sinh một lần; công thức~\eqref{eq:rag} vẫn
là cách hiểu đúng: chất lượng câu trả lời bị chặn trên bởi chất lượng truy xuất. Nếu đoạn văn chứa đáp án không
nằm trong top-$k$, bộ sinh dù mạnh tới đâu cũng chỉ còn cách đoán.

Hình~\ref{fig:bonho} đặt RAG trong khung hai loại bộ nhớ. \textbf{Bộ nhớ tham số} là trọng số LLM: rộng, nén, học
một lần, không sửa được từng mẩu và không chỉ ra được nguồn. \textbf{Bộ nhớ phi tham số} là kho tri thức ngoài:
cập nhật bằng cách thêm hay sửa tài liệu, mỗi câu trả lời trỏ được về đoạn văn cụ thể. RAG nối hai bộ nhớ qua bộ
truy xuất.

\hinhtikz{t1_memory}{Hai loại bộ nhớ của một hệ RAG: bộ nhớ tham số nằm trong trọng số LLM, bộ nhớ phi tham số là
kho tri thức ngoài; bộ truy xuất là cầu nối}{bonho}

\subsection{Quá trình phát triển}
Biểu diễn tri thức bằng đồ thị không mới. Mạng ngữ nghĩa có từ thập niên 1960; Web ngữ nghĩa của Berners-Lee năm
2001~\cite{bernerslee2001semantic} chuẩn hóa nó thành RDF, OWL và SPARQL; năm 2012 Google đưa thuật ngữ
``Knowledge Graph'' vào sản phẩm tìm kiếm. Nhánh thứ hai là mô hình ngôn ngữ: kiến trúc Transformer
năm 2017~\cite{vaswani2017attention} mở đường cho LLM, rồi RAG năm 2020 nối LLM với kho văn bản. Hai nhánh gặp
nhau từ 2023~\cite{pan2024unifying}: ReAct~\cite{yao2023react} cho LLM gọi công cụ trong vòng suy luận, GraphRAG năm
2024~\cite{edge2024graphrag} dùng đồ thị và cộng đồng để trả lời câu hỏi tổng hợp (Hình~\ref{fig:dongthoigian}).

\hinhtikz{t1_timeline}{Dòng thời gian của biểu diễn tri thức và RAG: từ mạng ngữ nghĩa, Web ngữ nghĩa, Google
Knowledge Graph tới Transformer, RAG, ReAct và GraphRAG}{dongthoigian}

\section{Vai trò và ý nghĩa của Knowledge Base và RAG}
Cơ sở tri thức (Knowledge Base, KB) là tập tri thức của một miền được tổ chức để máy truy cập được; trong bài này
KB có hai dạng song song: kho văn bản Markdown (dạng phi cấu trúc, phục vụ RAG) và hai KG (dạng có cấu trúc).
Giá trị của KB với một hệ thống dùng LLM nằm ở năm điểm.

\begin{enumerate}
  \item \textbf{Giảm ảo giác.} Ảo giác (hallucination) là khi mô hình sinh nội dung trôi chảy nhưng sai hoặc không
        có căn cứ~\cite{ji2023hallucination}. Ở Mục~\ref{sec:agent-ketqua}, cùng mô hình Qwen2.5-1.5B trả lời
        \SL{dem/agent_scen} câu hỏi kiểm thử: khi không có công cụ, câu trả lời chỉ chứa
        \SL{agent/agg/all/coverage_llm_onlyPct}\% dữ kiện bắt buộc; khi có KB qua tác tử, tỉ lệ này là
        \SL{agent/agg/all/coveragePct}\%.
  \item \textbf{Cập nhật không cần huấn luyện lại.} Nghị định 232/2025/NĐ-CP sửa Nghị định 24/2012 về vàng; với KB,
        chỉ cần thêm một thực thể và một cạnh \texttt{SUA\_DOI}, không phải tinh chỉnh mô hình.
  \item \textbf{Truy vết nguồn.} Trong y tế, một khuyến nghị không có nguồn thì không dùng được. Mọi cạnh trong
        hai KG đều mang mã tài liệu nguồn (\SL{kg/health/with_source} trên \SL{kg/health/edges} cạnh y tế,
        \SL{kg/finance/with_source} trên \SL{kg/finance/edges} cạnh tài chính).
  \item \textbf{Bảo mật dữ liệu.} Hồ sơ xét nghiệm của người bệnh không rời máy: LLM, embedder, reranker và cả hai
        KG chạy cục bộ.
  \item \textbf{Suy luận có kiểm soát.} Quy tắc nghiệp vụ (chống chỉ định, dấu tác động vĩ mô) nằm trong KG dưới
        dạng tường minh; LLM chỉ diễn đạt lại, không phải tự suy ra.
\end{enumerate}

Bảng~\ref{tab:sosanh-kb} so sánh ba cách đưa tri thức miền vào một hệ thống dùng LLM.

\begin{table}[H]\centering
\caption{So sánh LLM đơn lẻ, tinh chỉnh (fine-tuning) và hệ thống tích hợp KB + RAG}\label{tab:sosanh-kb}
{\small\begin{tabular}{L{3.3cm}L{3.5cm}L{3.5cm}L{3.9cm}}
\toprule
Tiêu chí & LLM đơn lẻ & Tinh chỉnh & KB + RAG (+ KG) \\
\midrule
Nguồn tri thức & tham số, cố định lúc huấn luyện & tham số, học thêm dữ liệu miền & kho ngoài, đọc lúc trả lời \\
Cập nhật & không & huấn luyện lại, tốn GPU & thêm, sửa tài liệu hoặc bộ ba \\
Truy vết nguồn & không & không & có, tới từng đoạn hoặc cạnh \\
Ảo giác & cao với tri thức hiếm & giảm nhưng vẫn còn & thấp nếu có bộ kiểm chứng \\
Câu hỏi đa bước & tự suy, khó kiểm soát & tự suy & duyệt đường trên KG \\
Chi phí triển khai & thấp & cao & trung bình: chỉ mục + truy xuất \\
Độ trễ & một lượt sinh & một lượt sinh & thêm truy xuất, rerank \\
Phù hợp & hội thoại chung & đổi văn phong, định dạng & y tế, pháp lý, tài chính \\
\bottomrule
\end{tabular}}
\end{table}

Tinh chỉnh và RAG không loại trừ nhau: tinh chỉnh dạy mô hình \emph{cách} trả lời (văn phong, định dạng trích dẫn),
RAG cung cấp \emph{nội dung} để trả lời. Với một mô hình nhỏ như Qwen2.5-1.5B, Chương~\ref{chap:agent} cho thấy mô
hình hiểu nội dung bằng chứng tốt nhưng thường quên ghi số hiệu trích dẫn; đó chính là phần việc mà tinh chỉnh
có thể bù.

\section{Bài toán và cách tiếp cận của tiểu luận}
Đề bài đặt hai bài toán ứng dụng có tính chất rất khác nhau, và chính sự khác nhau đó làm rõ khi nào cần KG,
khi nào cần RAG, khi nào cần mô hình học sâu.

\begin{itemize}
  \item \textbf{Y tế} là miền của \emph{ràng buộc}: tiêu chuẩn chẩn đoán có ngưỡng cụ thể, thuốc có chống chỉ định,
        thực phẩm có chỉ số đường huyết. Câu trả lời phải đúng từng con số và có nguồn. KG lưu ràng buộc, RAG cung
        cấp lời giải thích, MLP ước lượng nguy cơ từ hồ sơ xét nghiệm.
  \item \textbf{Tài chính} là miền của \emph{quan hệ nhân quả và bất định}: tỷ giá tác động lên cổ phiếu qua nhiều
        kênh, có kênh trái dấu nhau; giá tương lai không biết trước. KG lưu chiều tác động, dữ liệu thị trường kiểm
        chứng lại chiều đó, BiLSTM dự báo kèm khoảng bất định.
\end{itemize}

Cả hai miền dùng chung một lõi: lớp \texttt{KnowledgeGraph} với bản thể học, cùng một kho văn bản và đường ống
truy xuất, cùng một LLM cục bộ và cùng một tác tử ReAct. Mỗi thành phần được đánh giá riêng trước khi ghép:
truy xuất ở Chương~\ref{chap:rag}, KG ở Chương~\ref{chap:kg}, \ref{chap:yte}, \ref{chap:taichinh}, mô hình học
sâu ở Chương~\ref{chap:yte}, \ref{chap:taichinh}, và toàn hệ thống ở Chương~\ref{chap:agent}.
